
\documentclass{article}
\usepackage{colortbl}
\usepackage{makecell}
\usepackage{multirow}
\usepackage{supertabular}

\begin{document}

\newcounter{utterance}

\centering \large Interaction Transcript for game `cladder', experiment `full\_v1.5\_default', episode 3821 with Qwen3.5{-}9B.
\vspace{24pt}

{ \footnotesize  \setcounter{utterance}{1}
\setlength{\tabcolsep}{0pt}
\begin{supertabular}{c@{$\;$}|p{.15\linewidth}@{}p{.15\linewidth}p{.15\linewidth}p{.15\linewidth}p{.15\linewidth}p{.15\linewidth}}
    \# & \multicolumn{2}{c}{Player} && \multicolumn{2}{c}{Game Master} \\
    \hline

    \theutterance \stepcounter{utterance}  
    & & & \multicolumn{4}{p{0.6\linewidth}}{
        \cellcolor[rgb]{0.9,0.9,0.9}{
            \makecell[{{p{\linewidth}}}]{
                \texttt{\tiny{[P1$\langle$GM]}}
                \texttt{You are an expert in causal inference. The following question is not a typical commonsense query, but rather a meticulously designed question created by a professor specializing in causal inference, intended to assess the students' mastery of the course content.} \\
\texttt{Imagine a self{-}contained, hypothetical world with only the following conditions, and without any unmentioned factors or causal relationships: Gender has a direct effect on department competitiveness and peanut allergy. Department competitiveness has a direct effect on peanut allergy. For individuals who are not male, the probability of being allergic to peanuts is 48\%. For individuals who are male, the probability of being allergic to peanuts is 36\%. For individuals who are male, would it be more likely to see being allergic to peanuts if the individual had not been male?} \\
\texttt{Start your answer with 'Yes' or 'No', followed by additional reasoning or evidence to support your explanation.} \\
            }
        }
    }
     \\ \\

    \theutterance \stepcounter{utterance}  
    & \multicolumn{4}{p{0.6\linewidth}}{
        \cellcolor[rgb]{0.561,0.737,0.561}{
            \makecell[{{p{\linewidth}}}]{
                \texttt{\tiny{[P1$\rangle$GM]}}
                \texttt{No} \\
\\ 
\texttt{**Reasoning:**} \\
\\ 
\texttt{To determine the answer, we must analyze the causal structure and the specific probabilities provided in the hypothetical world.} \\
\\ 
\texttt{1.  **Identify the Causal Structure:**} \\
\texttt{    The problem states:} \\
\texttt{    *   Gender \$\textbackslash{}rightarrow\$ Peanut Allergy (Direct effect)} \\
\texttt{    *   Department Competitiveness \$\textbackslash{}rightarrow\$ Peanut Allergy (Direct effect)} \\
\texttt{    *   Gender \$\textbackslash{}rightarrow\$ Department Competitiveness (Direct effect)} \\
\\ 
\texttt{    This forms a chain structure (or a "mediator" structure) where Gender influences Department Competitiveness, which in turn influences Peanut Allergy. However, the question asks about the counterfactual for a **male** individual regarding their allergy status if they had **not** been male.} \\
\\ 
\texttt{2.  **Analyze the Counterfactual:**} \\
\texttt{    The question asks: "For individuals who are male, would it be more likely to see being allergic to peanuts if the individual had not been male?"} \\
\\ 
\texttt{    Mathematically, this compares two probabilities for a male individual (\$M=1\$):} \\
\texttt{    *   \$P(\textbackslash{}text\{Allergy\} | M=1)\$: The observed probability of being allergic given the individual is male.} \\
\texttt{    *   \$P(\textbackslash{}text\{Allergy\} | M=0)\$: The counterfactual probability of being allergic if the individual had been female (not male).} \\
\\ 
\texttt{3.  **Extract the Data:**} \\
\texttt{    *} \\
            }
        }
    }
    & & \\ \\

    \theutterance \stepcounter{utterance}  
    & & & \multicolumn{2}{p{0.3\linewidth}}{
        \cellcolor[rgb]{0.9,0.9,0.9}{
            \makecell[{{p{\linewidth}}}]{
                \texttt{\tiny{[GM$|$GM]}}
                \texttt{No} \\
            }
        }
    }
    & & \\ \\

    \theutterance \stepcounter{utterance}  
    & & & \multicolumn{2}{p{0.3\linewidth}}{
        \cellcolor[rgb]{0.9,0.9,0.9}{
            \makecell[{{p{\linewidth}}}]{
                \texttt{\tiny{[GM$|$GM]}}
                \texttt{yes} \\
            }
        }
    }
    & & \\ \\

    \theutterance \stepcounter{utterance}  
    & & & \multicolumn{2}{p{0.3\linewidth}}{
        \cellcolor[rgb]{0.9,0.9,0.9}{
            \makecell[{{p{\linewidth}}}]{
                \texttt{\tiny{[GM$|$GM]}}
                \texttt{game\_result = LOSE} \\
            }
        }
    }
    & & \\ \\

\end{supertabular}
}

\end{document}
